% Generated draft: requires mathematical and source review.
\documentclass[11pt]{article}
\usepackage[margin=25mm]{geometry}
\usepackage{fontspec}
\setmainfont{Latin Modern Roman}
\usepackage{amsmath,amssymb,mathtools,unicode-math}
\setmathfont{Latin Modern Math}
\usepackage{enumitem,needspace,xurl,hyperref}
\setlength{\parindent}{0pt}
\setlength{\parskip}{5pt}
\setlength{\emergencystretch}{3em}
\newcommand{\N}{\mathbb{N}}
\newcommand{\Z}{\mathbb{Z}}
\newcommand{\Q}{\mathbb{Q}}
\newcommand{\R}{\mathbb{R}}
\newcommand{\C}{\mathbb{C}}
\newcommand{\F}{\mathbb{F}}
\newcommand{\A}{\mathbb{A}}
\newcommand{\PP}{\mathbb{P}}
\newcommand{\E}{\mathbb{E}}
\newcommand{\eps}{\varepsilon}
\newcommand{\dd}{\,\mathrm{d}}
\begin{document}
\section*{UnsolvedMath ID: 30004074}
\textbf{Source problem code:} catalogue entry
\subsection{Definitions and mathematical statement}
Let \(d\in\N\) with \(d\geq5\). In the usual real-valued normalization, the focusing cubic wave equation is
\[
\partial_t^2u-\Delta_xu=u^3.
\]
Here \(\Delta_x=\sum_{j=1}^d\partial_{x_j}^2\). This normalization explains the conventional equation name; the supplied record does not specify the spatial domain, boundary conditions, solution regularity, initial-data space, or whether radial symmetry is assumed.

For terminology, the usual backward self-similar ansatz about a time \(T\) and a spatial point \(x_0\) is
\[
u(t,x)=\frac{1}{T-t}\,f\!\left(\frac{x-x_0}{T-t}\right),\qquad t<T,
\]
on a region where this expression is defined. The symbol \(f\) denotes a time-independent profile. In a radial formulation, one instead writes the profile as a function of \(\lvert x-x_0\rvert/(T-t)\). The record does not specify which formulation is intended, the profile domain, or a formula or characterizing conditions for its distinguished profile \(f^*\).

\emph{Source question, with unresolved qualifications.} For the focusing cubic wave equation in dimensions \(d\geq5\), does the source-designated profile \(f^*\) describe solutions at the threshold between stable blowup and global forward existence, and also describe solutions near that threshold during at least some intermediate times?

Here global forward existence means existence for all future times in the intended solution class. A precise meaning of stable blowup requires a topology on initial data and a specified blowup behavior that persists under sufficiently small perturbations. A precise threshold likewise requires the intended initial-data space and a definition of the separating set. The source does not supply these choices. It also does not define the rescaling, norm, convergence or approximation criterion, time regime, or quantifiers intended by ``describe'' and ``at least some intermediate times.'' Consequently, the question above preserves the supplied target but cannot be promoted to a fully quantified conjecture without additional source information.
\subsection{Short English statement}
The question asks whether the distinguished self-similar profile \(f^*\) governs solutions at the boundary between stable blowup and global existence, and near that boundary during at least some intermediate times, for \(d\geq5\). The supplied record leaves \(f^*\), the solution setting, and the required approximation behavior undefined. Its open status is supplied dated curation.
\subsection{Sources}
\begin{itemize}
\item Supplied local UnsolvedMath registry record, permanent ID \(30004074\), title \emph{Critical Self-Similar Profiles for Focusing Cubic Wave Blowup}, source problem code \texttt{catalogue entry}, collection \emph{Oberwolfach Reports Open Problems}. The proposed year is recorded as 2019; no proposer is identified. The registry labels the statement \texttt{corrected\_verified}; this is supplied metadata and does not constitute independent mathematical verification of this draft.
\item Supplied report citation: \emph{Nonlinear Evolution Equations: Analysis and Numerics} (2020), Oberwolfach Reports, pp.\ 305--405, DOI \texttt{10.4171/owr/2019/5}. \url{https://doi.org/10.4171/owr/2019/5}. The record describes Oberwolfach Reports as documentation of mathematical workshops and their open-problem discussions. No problem-specific page or passage is supplied.
\item Supplied literature source, labeled Glogic--Schoerkhuber: \url{https://doi.org/10.1016/j.aim.2021.107930}.
\item Supplied literature source, labeled Glogic--Maliborski--Schoerkhuber: \url{https://doi.org/10.1088/1361-6544/ab6f4d}.
\item Literature assessment dated exactly \texttt{2026-08-22}: the supplied curation labels the problem open and reports that primary work constructs the nontrivial self-similar profile and proves stability in some cases, while its role as the complete threshold between blowup and global existence remains analytically open. This is supplied dated curation derived from the OWR record, not independently verified. Its time-sensitive status and cited claims require independent checking.
\end{itemize}
\end{document}
